\section{Encoding a formula and its evaluation domain}
\label{sec:encoding}

We now turn the conditions imposed by a rational formula into finitely many
affine-pencil invertibility conditions.  We keep the original formula tree:
an inverse gate must have an invertible operand even when its contribution
later cancels or is multiplied by zero.
Affine-matrix representations of free-field elements appear in
Cohn and Reutenauer~\citep[Section~3.1]{CR94}.
Linear realizations by block inverses and tests for the original domain
appear in Hrube\v{s} and Wigderson~\citep[Theorem~2.6, Section~6,
and Proposition~7.1]{HW15}. Their domain criterion retains pencils
for all gate subformulas. We give the recursions and the domain
argument in full for the precise tree-size convention used here,
and show why the inverse-node pencils together with one root
augmentation suffice. The argument concerns the given tree and
requires no domain-preserving balancing transformation.

For a subtree $h$, let $|h|$ denote its number of vertices, let $\lambda(h)$
denote its number of leaves, and let $\iota(h)$ denote its number of inverse
gates.  Each occurrence of a subtree is treated separately.  A pencil of
order $q$ has the form
\[
 L(x)=A_0+\sum_{i=1}^n A_i x_i,
 \qquad A_0,\ldots,A_n\in\Mat_q(\Q).
\]
At a tuple $X=(X_1,\ldots,X_n)$ of matrices of order $d$ over a field
extension of $\Q$, its evaluation is
\[
 L(X)=A_0\otimes I_d+\sum_{i=1}^n A_i\otimes X_i.
\]
When a constant row $u$ or column $v$ is used with such an evaluation, we
write $u$ and $v$ for $u\otimes I_d$ and $v\otimes I_d$, respectively.
Thus the expression $uL(X)^{-1}v$ is a matrix of order $d$.

\begin{lemma}[Affine realization of a subtree]
\label{lem:formula-realization}
For each subtree $h$ there are an affine pencil $L_h$ of order $q_h$, a
constant rational row $u_h$ of length $q_h$, and a constant rational column
$v_h$ of length $q_h$, such that
\begin{equation}
 q_h=2\lambda(h)+\iota(h)=|h|+1.
 \label{eq:realization-order}
\end{equation}
At every matrix tuple over every field extension of $\Q$ where the
original subtree $h$ is defined, $L_h$ is invertible and
\begin{equation}
 h=u_hL_h^{-1}v_h.
 \label{eq:realization-invariant}
\end{equation}
\end{lemma}

\begin{proof}
Here are the constructions, followed by verification of the invariant.
All zero blocks have the sizes prescribed by the adjacent blocks.

\emph{Leaf.}  If $h$ is a variable or a rational constant, set
\[
 L_h=\begin{pmatrix}1&-h\\0&1\end{pmatrix},
 \qquad u_h=\begin{pmatrix}1&0\end{pmatrix},
 \qquad v_h=\begin{pmatrix}0\\1\end{pmatrix}.
\]

\emph{Addition.}  For $h=f+g$, set
\[
 L_h=\begin{pmatrix}L_f&0\\0&L_g\end{pmatrix},
 \qquad u_h=\begin{pmatrix}u_f&u_g\end{pmatrix},
 \qquad v_h=\begin{pmatrix}v_f\\v_g\end{pmatrix}.
\]

\emph{Ordered multiplication.}  For $h=fg$, set
\[
 L_h=\begin{pmatrix}L_f&-v_fu_g\\0&L_g\end{pmatrix},
 \qquad u_h=\begin{pmatrix}u_f&0\end{pmatrix},
 \qquad v_h=\begin{pmatrix}0\\v_g\end{pmatrix}.
\]

\emph{Inversion.}  For $h=g^{-1}$, set
\[
 L_h=\begin{pmatrix}L_g&v_g\\u_g&0\end{pmatrix},
 \qquad u_h=\begin{pmatrix}0&-1\end{pmatrix},
 \qquad v_h=\begin{pmatrix}0\\1\end{pmatrix},
\]
where the zero part of each selector has length $q_g$.  Equivalently,
$u_h=-e_{q_g+1}^{\mathsf T}$ and $v_h=e_{q_g+1}$.

The rows and columns just specified are constant.  In particular,
$v_fu_g$ is a constant matrix, so every displayed pencil is affine in
the original variables.  No multiplication of variable entries is
introduced into a pencil.

Fix an evaluation tuple and prove the invariant by induction on the
subtree.  A leaf pencil has inverse
$\left(\begin{smallmatrix}I_d&h\\0&I_d\end{smallmatrix}\right)$, and its selectors
give $h$.  At an addition gate, definedness of the gate implies
definedness of both children.  Their pencils are invertible by induction,
and block-diagonal inversion gives the sum of the two represented values.
At a multiplication gate, the same induction gives
\[
 L_h^{-1}=
 \begin{pmatrix}
  L_f^{-1}&L_f^{-1}v_fu_gL_g^{-1}\\
  0&L_g^{-1}
 \end{pmatrix}.
\]
The selected block is consequently
\[
 u_fL_f^{-1}v_fu_gL_g^{-1}v_g
   =(u_fL_f^{-1}v_f)(u_gL_g^{-1}v_g)=fg,
\]
with the required order of multiplication.

At an inverse gate $h=g^{-1}$, definedness means that $g$ is defined and
its value is invertible.  By induction $L_g$ is invertible and
$u_gL_g^{-1}v_g=g$.  The block factorization
\begin{equation}
 \begin{pmatrix}L_g&v_g\\u_g&0\end{pmatrix}
 =
 \begin{pmatrix}I&0\\u_gL_g^{-1}&I\end{pmatrix}
 \begin{pmatrix}L_g&0\\0&-g\end{pmatrix}
 \begin{pmatrix}I&L_g^{-1}v_g\\0&I\end{pmatrix}
 \label{eq:inverse-schur-factorization}
\end{equation}
shows that $L_h$ is invertible.  The lower-right block of its inverse is
$(-g)^{-1}=-g^{-1}$, so the negative row selector gives precisely
$g^{-1}$.  This proves the invariant in all four cases.

The orders are $2$ at leaves, add at binary gates, and increase by one
at inverse gates.  Hence $q_h=2\lambda(h)+\iota(h)$.  A rooted tree whose
internal vertices have one or two children has $\lambda(h)-1$ binary
vertices: counting edges both by vertices and by children gives this
identity.  Thus
\[
 |h|=\lambda(h)+(\lambda(h)-1)+\iota(h)
     =2\lambda(h)+\iota(h)-1,
\]
which proves~\eqref{eq:realization-order}.
\end{proof}

For the root formula $\Phi$, define its \emph{root augmentation} by
\begin{equation}
 J_\Phi=
 \begin{pmatrix}L_\Phi&v_\Phi\\u_\Phi&0\end{pmatrix}.
 \label{eq:root-augmentation}
\end{equation}
This has the same block form as an inverse-node pencil, but no inverse
gate is being added to the formula or counted in its size.

\begin{proposition}[Pencil certificates for the original formula tree]
\label{prop:formula-certificates}
Let $\Phi$ be a formula of size at most $s$, and put $B=2s+1$.  Form the
indexed collection $\mathcal C_\Phi$ consisting of $L_h$ for every
inverse-node occurrence $h$ of $\Phi$, together with $J_\Phi$.
Then:
\begin{enumerate}
 \item The collection contains at most $s+1$ affine pencils over $\Q$,
       each of positive order at most $B$.
 \item At every tuple of matrices of a common positive order over a
       field extension of $\Q$, all pencils in $\mathcal C_\Phi$ are
       invertible if and only if the original formula tree $\Phi$ is
       defined there and its value is invertible.
 \item If $\Phi$ is admissible and nonzero, there is a single tuple over
       some field extension of $\Q$ at which every pencil in
       $\mathcal C_\Phi$ is invertible.
\end{enumerate}
Adding the order-one constant pencil $1$ for normalization gives at
most $s+2$ pencils, still of order at most $B$, with the same simultaneous
feasibility property in the third assertion.
\end{proposition}

\begin{proof}
There are at most $s$ inverse-node occurrences.  By
Lemma~\ref{lem:formula-realization}, each corresponding pencil has order
$|h|+1\le s+1\le B$.  The augmentation has order
$|\Phi|+2\le s+2\le2s+1=B$, where the last inequality uses $s\ge1$.
This proves the first assertion.  The extra constant pencil has order
one and evaluates to an identity matrix at every tuple.

For one implication in the second assertion, assume that
$\Phi$ is defined and its value is invertible.  Each subtree is then
defined, so Lemma~\ref{lem:formula-realization} makes every inverse-node
pencil invertible.  The same lemma makes $L_\Phi$ invertible and gives
$u_\Phi L_\Phi^{-1}v_\Phi=\Phi$.  Applying
factorization~\eqref{eq:inverse-schur-factorization} at the root shows
that $J_\Phi$ is invertible, because its Schur complement is $-\Phi$.

Conversely, suppose all pencils in $\mathcal C_\Phi$ are invertible.
First establish that the original tree is defined, proceeding from
the leaves toward the root.  Leaves are always defined.  At an addition
or multiplication gate, both children have already been shown defined,
so the gate is defined.  At an inverse gate $h=g^{-1}$, all gates in
the child subtree have already been handled; in particular $g$ is
defined.  Lemma~\ref{lem:formula-realization} therefore makes $L_g$
invertible and represents $g$ by $u_gL_g^{-1}v_g$.  The pencil $L_h$ is
one of the assumed invertible pencils.  The two outer factors in
\eqref{eq:inverse-schur-factorization} are invertible, so the middle
factor is invertible as well.  It follows that $g$ is invertible, which
is exactly the condition needed for this inverse gate to be defined.

This induction proves definedness of the entire original tree.  It
also explains why no separate constraint on the pencil of an ordinary
subtree is needed: whenever that pencil is used, its invertibility
already follows from the established definedness of its subtree.
Now $L_\Phi$ is invertible by the realization lemma.  The assumed
invertibility of $J_\Phi$, whose Schur complement is $-\Phi$, forces
the final value $\Phi$ to be invertible.  This proves the reverse
implication without simplifying the formula or enlarging its domain.

Finally, if $\Phi$ is admissible and nonzero,
Lemma~\ref{lem:invertible-witness} supplies an evaluation over a field
extension of $\Q$ at which the original tree is defined and its value
is invertible.  The second assertion makes every pencil in
$\mathcal C_\Phi$ invertible at this one evaluation.  The constant
pencil $1$ is invertible there as well.  This proves the third assertion
and the normalization statement.
\end{proof}

\begin{remark}[Why every inverse occurrence is retained]
Consider the scalar expression
\[
 a=x^{-1},\qquad b=a^{-1},\qquad \Phi=b+1.
\]
It is undefined at $x=0$, although its simplified rational function is
$x+1$.  For the pencils above, direct computation gives
\[
 \det L_a=-x,\qquad \det L_b=1,\qquad \det J_\Phi=-(x+1).
\]
Thus the outer inverse pencil and the root augmentation alone would
accept $x=0$.  The retained inner inverse pencil $L_a$ rejects it.
Likewise, the formula $1+0\cdot x^{-1}$ retains the requirement that
$x$ be invertible.  The certificate collection is indexed by gates of
the given tree, including gates in branches whose later contributions
cancel.
\end{remark}

The witnesses and pencils in this section are used to prove the hitting
property; the generator will receive only $n,s$ and will not construct
them from an unknown formula.  The order and count bounds depend only
on the number of vertices.  They impose no restriction on inverse
nesting or on the bit lengths of the rational leaf constants.
